\documentclass[11pt]{article}
\usepackage[margin=1in]{geometry}
\usepackage{amsmath,amssymb,amsthm}
\usepackage{booktabs}
\usepackage{float}
\usepackage{algorithm}
\usepackage{algpseudocode}
\algnotext{EndFor}
\usepackage{tikz}
\usetikzlibrary{arrows.meta}
\usepackage[colorlinks=true,linkcolor=blue,urlcolor=blue]{hyperref}

\newtheorem{bound}{Bound}
\theoremstyle{definition}
\newtheorem*{claim}{Claim}

\title{CS 422 Assignment 4: Maximum Concurrent Flow and Network Design}
\author{Group 6: Vraj Patel, Rahul Siddharth, Priyanka Soe, Jay Sunkara}
\date{}

\begin{document}
\maketitle
\vspace{-3em}

\begin{center}
Code repository: \url{https://github.com/vrajpatel-18/CS422-hw4}
\end{center}

\section{Part 1: Uniform Traffic Matrix}\label{p1}

\subsection{Answer}

Number the nodes $0,\dots,n-1$ and let $k = (n-1)/d$. Every node uses the same $d$ \emph{jumps}: node $i$ links to
$(i+j) \bmod n$ for each jump $j$, where the jumps are
\[
J = \{\, x \in \{1,\dots,n-1\} \;:\; x^d \equiv 1 \pmod n \,\}.
\]
For $n$ prime, $J$ has exactly $d$ elements (this is where $d \mid n-1$ is used). No topology can beat
\[
\theta^* \;=\; \frac{k}{2k-1},
\]
and our topology reaches it whenever every node is within 2 hops of every other. For example, $n=13$, $d=4$ gives
$J=\{1,5,8,12\}$, every node is within 2 hops (checked in Section~\ref{p1:why}), and $\theta = \theta^* = 3/5$.

\paragraph{Scope.} The upper bound holds for every $n$ and every topology. Finding the optimal topology for every $n$
and $d$ is closely tied to the degree--diameter problem (how many nodes a degree-$d$ graph can fit within a given
number of hops), which is a long-standing open problem. So the answer here is the bound, plus a construction that
meets it exactly for prime $n$ whenever it has every node within 2 hops (Section~\ref{p1:lim}).

\subsection{Setup}

Each node has $d$ outgoing and $d$ incoming links of capacity 1, so the network has $nd$ units of capacity in total.
Each node sends $\frac{d}{n-1}$ to each of the other $n-1$ nodes. The \emph{concurrent flow} $\theta$ is the largest
factor for which $\theta T$ fits on the links. Notation:
\begin{itemize}
\item $\mathrm{hops}(s,t)$: the fewest links needed to get from $s$ to $t$.
\item $A$: the average of $\mathrm{hops}(s,t)$ over all $n(n-1)$ ordered pairs.
\item $k = (n-1)/d$: the number of destinations per outgoing link (an integer by assumption).
\end{itemize}

\subsection{Upper bounds that hold for every topology}

\begin{bound}\label{b:1}
For every topology, $\theta \le 1/A$.
\end{bound}
\begin{proof}
Each unit of traffic from $s$ to $t$ crosses at least $\mathrm{hops}(s,t)$ links and uses one unit of capacity on
each. So the scaled traffic uses at least
\[
\theta\cdot\frac{d}{n-1}\cdot\sum_{s\ne t}\mathrm{hops}(s,t) \;=\; \theta \cdot nd \cdot A
\]
units of capacity. Only $nd$ units exist, so $\theta\, nd\, A \le nd$, i.e.\ $\theta \le 1/A$.
\end{proof}

\noindent Bound~\ref{b:1} is met with equality only if (i) all traffic follows shortest routes and (ii) every link is
completely full. So a good topology needs \textbf{few hops} and \textbf{evenly spread load}.

\begin{bound}\label{b:2}
For every topology, $A \ge \frac{2k-1}{k}$, so $\theta \le \frac{k}{2k-1}$.
\end{bound}
\begin{proof}
A node has only $d$ outgoing links, so at most $d$ of its $n-1$ destinations, a fraction $1/k$, are 1 hop away. All
others are at least 2 hops away. So $A \ge \frac{1}{k}\cdot 1 + \frac{k-1}{k}\cdot 2 = \frac{2k-1}{k}$, and
Bound~\ref{b:1} gives $\theta \le \frac{k}{2k-1}$.
\end{proof}

\begin{table}[H]
\centering
\begin{tabular}{l ccccc}
\toprule
$k = (n-1)/d$ & 1 & 2 & 3 & 4 & large \\
\midrule
best possible $\theta^*$ & 1 & 2/3 & 3/5 & 4/7 & $\to 1/2$ \\
\bottomrule
\end{tabular}
\end{table}
\noindent The case $k=1$ ($d = n-1$) is the complete graph: every pair has a direct link and $\theta^*=1$.

\subsection{The construction}

\begin{algorithm}[H]
\caption{\textsc{BuildTopology}$(n,d)$, for $n$ prime and $d \mid n-1$}
\begin{algorithmic}[1]
\State $J \gets \{\, x \in \{1,\dots,n-1\} : x^d \bmod n = 1 \,\}$ \Comment{exactly $d$ jumps}
\For{each node $i = 0,\dots,n-1$}
  \For{each jump $j \in J$}
    \State add link $i \to (i+j) \bmod n$ \Comment{capacity 1}
  \EndFor
\EndFor
\State \textbf{Routing:} split the traffic of each pair evenly over all its shortest routes.
\end{algorithmic}
\end{algorithm}

\paragraph{Why it is a valid topology.} Node $i$ has $d$ outgoing links (one per jump) and $d$ incoming links (from
$i-j$ for each jump $j$). Jumps are distinct and nonzero, so there are no self-loops and no duplicate links. For $n$
prime, $x^d \equiv 1$ has exactly $d$ solutions when $d \mid n-1$ (the nonzero numbers mod a prime form a cyclic group
of size $n-1$), so $|J| = d$.

\paragraph{The key property.} The jumps are \emph{closed under multiplication}: if $x^d \equiv 1$ and $y^d \equiv 1$ then
$(xy)^d \equiv 1$. Likewise, the inverse of a jump is a jump. For $n=7$, $d=3$: $J=\{1,2,4\}$, and
$2\cdot 2=4$, $2\cdot4=8\equiv1$, $4\cdot4=16\equiv2$. Finding $J$ by hand is easy: $1^3=1$, $2^3=8\equiv1$,
$4^3=64\equiv1 \pmod 7$, but $3^3 = 27 \equiv 6$, so $3 \notin J$.

\begin{figure}[ht]
\centering
\tikzset{
  v/.style={circle, draw, minimum size=6mm, inner sep=0pt, font=\small, fill=white},
  e/.style={-{Stealth[length=2mm]}, thick, shorten >=0.5pt},
  j1/.style={e, blue!70!black},
  j2/.style={e, red!75!black},
  j4/.style={e, green!50!black},
}
\begin{minipage}[b]{0.46\linewidth}
\centering
\begin{tikzpicture}
  \foreach \i in {0,...,6} \node[v] (a\i) at ({90-360*\i/7}:1.9) {\i};
  \foreach \i in {0,...,6} {
    \pgfmathtruncatemacro{\p}{mod(\i+1,7)}
    \pgfmathtruncatemacro{\q}{mod(\i+2,7)}
    \pgfmathtruncatemacro{\r}{mod(\i+4,7)}
    \draw[j1] (a\i) -- (a\p);
    \draw[j2] (a\i) -- (a\q);
    \draw[j4] (a\i) -- (a\r);
  }
  \foreach \c/\s/\x in {j1/1/-2.7, j2/2/-0.85, j4/4/1.0}
    \draw[\c, -] (\x,-2.6) -- ++(0.35,0) node[right, font=\footnotesize, text=black] {$i \to i{+}\s$};
\end{tikzpicture}\\[2pt]
{\small (a) the topology: jumps $\{1,2,4\}$}
\end{minipage}\hfill
\begin{minipage}[b]{0.46\linewidth}
\centering
\begin{tikzpicture}
  \node[v] (b0) at (0,0) {0};
  \node[v] (b1) at (2,1.2) {1};
  \node[v] (b2) at (2,0) {2};
  \node[v] (b4) at (2,-1.2) {4};
  \node[v] (b3) at (4,1.2) {3};
  \node[v] (b5) at (4,0) {5};
  \node[v] (b6) at (4,-1.2) {6};
  \node[font=\small] at (2,1.9) {1 hop};
  \node[font=\small] at (4,1.9) {2 hops};
  \draw[j1] (b0) -- (b1);  \draw[j2] (b0) -- (b2);  \draw[j4] (b0) -- (b4);
  \draw[j2] (b1) -- (b3);  \draw[j4] (b1) -- (b5);
  \draw[j1] (b2) -- (b3);  \draw[j4] (b2) -- (b6);
  \draw[j1] (b4) -- (b5);  \draw[j2] (b4) -- (b6);
  \node at (0,-2.6) {};
\end{tikzpicture}\\[2pt]
{\small (b) shortest routes from node 0}
\end{minipage}
\caption{$n=7$, $d=3$, $J=\{1,2,4\}$. (a) The topology; colours mark the jump of each link. (b) From node 0, the
3 jumps are 1 hop away and the other 3 nodes are 2 hops away, each reached by exactly two shortest routes
($3=1+2=2+1$, $5=1+4=4+1$, $6=2+4=4+2$). So $A = 1.5$ and $\theta = 2/3$, which matches Bound~\ref{b:2}.}
\label{fig:n7}
\end{figure}

\subsection{Why it maximizes the concurrent flow}\label{p1:why}

This section shows the construction meets both bounds with equality. Claim~1 handles Bound~\ref{b:1} (the load is
perfectly even), and Claim~2 handles Bound~\ref{b:2} (the hop count is as small as possible).

\paragraph{What a relabeling is.} A relabeling renames every node, giving node $x$ the new name $\sigma(x)$. We use
relabelings of one form:
\[
\sigma(x) = a\cdot x + b \pmod n, \qquad a \in J,\ b \text{ any number in } \{0,\dots,n-1\}.
\]
In words: multiply every node's number by a jump $a$, then shift every number by $b$. Each node gets a different new
name, because $n$ is prime, so multiplying by $a \ne 0$ can be undone. We say a relabeling \emph{maps the topology
onto itself} if, after renaming, the set of links is exactly the same as before. When that happens, nothing about the
network can tell two links apart: if one link had some special feature, the relabeling would carry that feature to
another link.

\paragraph{A worked example ($n=7$, $J=\{1,2,4\}$).} Take $a=4$, $b=2$, so $\sigma(x) = 4x+2 \bmod 7$:
\begin{center}
\begin{tabular}{l|ccccccc}
old name $x$ & 0 & 1 & 2 & 3 & 4 & 5 & 6 \\
new name $\sigma(x)$ & 2 & 6 & 3 & 0 & 4 & 1 & 5
\end{tabular}
\end{center}
The link $0\to1$ (jump 1) becomes $2\to6$, which is a link (jump 4). The link $0\to4$ (jump 4) becomes $2\to4$, also a
link (jump 2). Every link maps to a link this way. This particular $\sigma$ also sends the link $3\to5$ onto the link
$0\to1$: $\sigma(3)=0$ and $\sigma(5)=1$.

\begin{claim}[1: every link carries the same load, so $\theta = 1/A$]
\end{claim}
\begin{proof}
We check three facts about the relabelings $\sigma(x) = ax+b$.
\begin{itemize}
\item \emph{Each one maps the topology onto itself.} The link $x \to x+j$ is renamed to $\sigma(x) \to \sigma(x) + aj$.
  Since $a$ and $j$ are both jumps, $aj$ is a jump (the key property), so the renamed link is a real link.
\item \emph{Any link can be moved onto any other.} To move the link $i \to i+j$ onto $0 \to 1$, use $a = j^{-1}$ (a
  jump, since inverses of jumps are jumps) and $b = -j^{-1} i$. Then $\sigma(i) = 0$ and $\sigma(i+j) = j^{-1}j = 1$.
  In the example above, $j=2$, $j^{-1}=4$ and $i=3$ give exactly $\sigma(x)=4x+2$. Any link can be moved to $0\to1$
  and back out again, so any link can be moved onto any other.
\item \emph{Nothing else changes either.} The traffic is the same amount between every pair, and the routing rule
  (split evenly over shortest routes) only looks at hop counts. Renaming nodes changes neither.
\end{itemize}
Now suppose link $e$ carries some load. Pick a relabeling that moves $e$ onto another link $e'$. After renaming, the
network, the traffic and the routing are all exactly as before, so $e'$ must carry the same load that $e$ did. Hence
all $nd$ links carry the same load. The total load is $\theta \cdot nd \cdot A$ (from the proof of Bound~\ref{b:1}),
so each link carries $\theta A$. Setting $\theta = 1/A$ fills every link exactly to capacity. That routing is
feasible, and Bound~\ref{b:1} says nothing larger is, so $\theta = 1/A$.
\end{proof}

\begin{claim}[2: within 2 hops, $A$ meets Bound~\ref{b:2}]
\end{claim}
\begin{proof}
The 1-hop destinations of a node are exactly its $d$ jumps, the most Bound~\ref{b:2} allows. If every other node is
a sum of two jumps, it is exactly 2 hops away, so $A = \frac{2k-1}{k}$ and $\theta = \frac{k}{2k-1}$, the best
possible. Every node sees the same picture as node 0 (rename with $a=1$, which just shifts every number by $b$), so
checking node 0 suffices.
\end{proof}

\noindent Putting the two claims together: whenever every node is within 2 hops, our topology achieves
$\theta^* = \frac{k}{2k-1}$, and no topology can do better.

\paragraph{Checking the 2-hop condition by hand.} By Claim~2, it is enough to check that from node 0, every node
that is not a jump is a sum of two jumps (mod $n$).
\begin{itemize}
\item $n=7$, $d=3$, $J=\{1,2,4\}$: the sums of two jumps are $\{1,2,3,4,5,6\}$ (mod 7), which cover all of $\{3,5,6\}$
  (Figure~\ref{fig:n7}). So $\theta = 2/3 = \theta^*$.
\item $n=13$, $d=4$, $J=\{1,5,8,12\}$. Ignoring $0$, the sums of two jumps mod 13 are $\{2,\allowbreak 3,\allowbreak 4,\allowbreak 6,\allowbreak 7,\allowbreak 9,\allowbreak 10,\allowbreak 11\}$.
  Together with $J$ this is all 12 other nodes. So $\theta = 3/5 = \theta^*$.
\end{itemize}

\paragraph{Why both properties matter ($n=13$, $d=4$).}
\begin{itemize}
\item Jumps $\{1,2,3,4\}$: from node 0, four nodes each are at 1, 2 and 3 hops, so $A = 2$ and $\theta \le 1/2$ by
  Bound~\ref{b:1}. Too many hops.
\item Jumps $\{1,2,6,9\}$: also within 2 hops (the sums of two jumps give $\{2,3,4,5,7,8,10,11,12\}$), with the same
  $A = 5/3$ as ours. But $\{1,2,6,9\}$ is not closed under multiplication ($2\cdot2=4 \notin J$), so the relabelings of
  Claim~1 do not map this graph onto itself. Nothing then forces every link to carry the same load, and Bound~\ref{b:1}
  is only met if every link is exactly full. Few hops alone are not enough.
\item Our jumps $\{1,5,8,12\}$: within 2 hops, and every link carries the same load (Claim~1), so
  $\theta = 3/5 = 0.600$, the best possible.
\end{itemize}

\subsection{Limitations}\label{p1:lim}
\begin{itemize}
\item \textbf{$n$ must be prime} for this construction. Only then is $|J| = d$ guaranteed, and only then does every
  relabeling in Claim~1 give each node a different name. Bounds~\ref{b:1} and~\ref{b:2} hold for every $n$.
\item \textbf{Very sparse networks.} A node reaches at most $d$ nodes in 1 hop and $d^2$ in 2 hops. If
  $d + d^2 < n-1$, not everyone can be within 2 hops, so no topology reaches Bound~\ref{b:2}. (Counting $d^3$ nodes at
  3 hops, and so on, gives a tighter bound.) Even when $d + d^2 = n-1$ exactly, every 2-hop route would have to land
  on a different node, and it is known that no such graph exists for $d \ge 2$. For example, $n=13$, $d=3$ has
  $d+d^2 = 12 = n-1$, so no topology reaches Bound~\ref{b:2} there.
\item \textbf{The jumps do not always give 2 hops.} Even when the counting above allows 2 hops, our jumps can need 3.
  For $n=17$, $d=4$ ($d+d^2 = 20 \ge 16$), the jumps are $\{1,4,13,16\}$, and the sums of two jumps (mod 17, ignoring
  $0$) are $\{2,3,5,8,9,12,14,15\}$. Together with $J$ that is only 12 of the 16 other nodes: $6, 7, 10$ and $11$ are
  3 hops away. Then $A = (4\cdot1 + 8\cdot2 + 4\cdot3)/16 = 2$, so our $\theta = 1/2$ falls short of $k/(2k-1) = 4/7$.
\item In both of these cases the load is still perfectly even ($\theta = 1/A$ exactly, by Claim~1), so the only loss
  comes from the extra hops.
\end{itemize}

\section{Part 2: Arbitrary Traffic Matrix}
We chose to do the linear integer programming formulation for Part 2. The initial problem for the project specifies designing a network topology that maximizes the concurrent flow for the given traffic matrix. We are given a traffic matrix with 8 nodes, 4 incoming links, and 4 outgoing links per node, with each link having a capacity of 1 unit. We are also given the hose model traffic set, which essentially gives us the following constraints:

\begin{itemize}
    \item No node sends or receives units to/from itself, specified by:
    \[
    T_{ii} = 0 \quad \forall i
    \]

    \item The sum of all the units node $i$ can send out must be less than or equal to 4:
    \[
    \sum_{j=1}^{8} T_{ij} \leq 4 \quad \forall i
    \]

    \item The sum of all the units node $j$ can receive must be less than or equal to 4:
    \[
    \sum_{i=1}^{8} T_{ij} \leq 4 \quad \forall j
    \]
\end{itemize}

In addition, no nodes can send negative traffic, so $T_{ij} \geq 0$. From this problem, we can derive our decision variables, objective function, and constraints as follows.

\subsection{Decision Variables}

In this problem, there are two major things we are concerned about: 1) where to direct the 4 incoming and 4 outgoing links (existence of links) and 2) the units each link sends (flow). To account for 1, we create the first decision variable:

\[
x_{ij} \in \{0,1\}
\]

Where $x_{ij}$ is 1 if there exists a directed link from $i$ to $j$, and 0 otherwise. We assume that there can be at most one directed link from $i$ to $j$.

For flow, we use:

\[
f_{ij}^{st} \geq 0
\]

Where $f_{ij}^{st}$ represents the units $s$ sends to $t$ that flow across the link $i \rightarrow j$.

We use this variable because traffic may not always travel directly from its source to its destination. Instead, it may have to pass through intermediate nodes, so we need to keep track of how much traffic from each source-destination pair travels across each link.

We also define $K$ as the set of source-destination pairs with positive traffic demand:

\[
K = \{(s,t) \mid s \neq t,\ T_{st} > 0\}
\]

Finally, we introduce another decision variable:

\[
\lambda \geq 0
\]

Where $\lambda$ represents the concurrent flow factor. Essentially, this tells us what fraction or multiple of the requested traffic demands the network can support simultaneously. For example, if $\lambda = 0.7$, the network can support 70\% of each requested traffic demand.

\subsection{Objective Function}

Our goal is to maximize the concurrent flow for the given traffic matrix. Since $\lambda$ represents the concurrent flow factor, our objective function is:

\[
\boxed{\max \lambda}
\]

By maximizing $\lambda$, we are finding the network topology that can support the largest common scaling factor of the traffic demands while still satisfying the link capacity and topology constraints.

\subsection{Constraints}

We have several constraints to ensure that the topology we design is valid and that traffic can actually be routed through the network.

\subsubsection{Outgoing Link Constraint}

Each node must have exactly 4 outgoing links. Since $x_{ij}$ represents whether a directed link exists from $i$ to $j$, we can specify this as:

\[
\sum_{j=1}^{8} x_{ij} = 4 \quad \forall i
\]

This ensures that every node has exactly 4 outgoing connections.

\subsubsection{Incoming Link Constraint}

Similarly, each node must have exactly 4 incoming links:

\[
\sum_{i=1}^{8} x_{ij} = 4 \quad \forall j
\]

This ensures that every node has exactly 4 incoming connections.

\subsubsection{No Self-Links}

Since nodes do not need to send traffic to themselves, we do not allow a directed link from a node to itself:

\[
x_{ii} = 0 \quad \forall i
\]

\subsubsection{Link Capacity Constraint}

Each directed link has a capacity of 1 unit. Therefore, the total traffic flowing through any link cannot exceed 1.

We specify this as:

\[
\sum_{(s,t)\in K} f_{ij}^{st} \leq x_{ij}
\quad \forall i \neq j
\]

This constraint accounts for both the capacity of a link and whether the link exists.

If $x_{ij} = 0$, then no traffic can flow through that link because it does not exist. If $x_{ij} = 1$, then the sum of all traffic flowing through the link can be at most 1 unit.

We sum over all source-destination pairs because multiple traffic demands may use the same link, and their combined flow must still satisfy the capacity constraint.

\subsubsection{Flow Conservation Constraint}

We also need to ensure that traffic is properly routed from its source to its destination. In particular, any intermediate node must send out the same amount of traffic that it receives for a given source-destination pair.

For every $(s,t) \in K$ and every node $v$, we use:

\[
\sum_{j=1}^{8} f_{vj}^{st}
-
\sum_{i=1}^{8} f_{iv}^{st}
=
\begin{cases}
\lambda T_{st} & \text{if } v=s \\
-\lambda T_{st} & \text{if } v=t \\
0 & \text{otherwise}
\end{cases}
\]

This gives us three different cases:

\begin{itemize}
    \item If $v=s$, then the node is the source. Its net outgoing traffic must equal $\lambda T_{st}$, since this is the amount of traffic we are trying to send.

    \item If $v=t$, then the node is the destination. Its net incoming traffic must equal $\lambda T_{st}$, since this is the amount of traffic it needs to receive.

    \item Otherwise, the node is an intermediate node. Its incoming and outgoing traffic must be equal because it is only forwarding traffic.
\end{itemize}

This ensures that traffic cannot be created or lost while traveling through the network.

Finally, we also specify:

\[
f_{ii}^{st} = 0
\quad \forall i,\ (s,t)\in K
\]

This prevents traffic from flowing through self-links.

Together with the nonnegative flow variables and binary topology variables, these constraints ensure that the resulting topology and routing are valid.

\subsection{Intuition Behind the Formulation}

The main idea behind our formulation is that we are optimizing both the topology and the routing of traffic at the same time.

The binary variables $x_{ij}$ determine which directed links should exist, while the continuous variables $f_{ij}^{st}$ determine how traffic should be routed through those links.

The incoming and outgoing link constraints ensure that each node has exactly 4 incoming and 4 outgoing connections. The capacity constraints ensure that no link carries more than 1 unit of traffic, and the flow conservation constraints ensure that traffic is properly routed from its source to its destination.

By maximizing $\lambda$, we are finding the topology that can support the greatest common scaling factor of all the traffic demands.

Since the topology variables are binary and the flow variables are continuous, this is a mixed-integer linear programming problem (MILP). We implemented this formulation using the Gurobi Python API.

Our formulation works for any given traffic matrix in the hose model because the traffic matrix is provided as an input to the program. Depending on how traffic is distributed, Gurobi can choose a different topology to maximize the concurrent flow.

Under our assumption that there is at most one directed link between any two nodes and that traffic can be split across multiple paths, the formulation considers all valid topologies and feasible routings. Therefore, when Gurobi finds an optimal solution, no other solution satisfying these constraints can achieve a higher concurrent flow factor for the given traffic matrix.

\subsection{Implementation}

We implemented the formulation in Python using the Gurobi optimization library.

The program takes an $8 \times 8$ traffic matrix as input and first identifies all source-destination pairs with positive traffic demand. It then creates the binary topology variables, continuous flow variables, and concurrent flow factor.

Next, we add the incoming and outgoing degree constraints, link capacity constraints, and flow conservation constraints described above.

Finally, we use Gurobi to maximize $\lambda$. After optimization, the program outputs the optimal concurrent flow factor and the directed links selected for the resulting topology.

We also included a check for the case where the traffic matrix contains no positive demands. In this case, there is no traffic to route, so the concurrent flow factor is unbounded.

\subsection{Results}

To test our implementation, we used three different traffic matrices: uniform, sparse, and nonuniform. All three satisfy the hose model constraints, with each node having a total incoming and outgoing traffic demand of 4 units.

The results are summarized below:

\begin{table}[h]
\centering
\begin{tabular}{lcc}
\toprule
\textbf{Traffic Matrix} & \textbf{Optimal $\lambda$} & \textbf{Solve Time (s)} \\
\midrule
Uniform & 0.7000 & 0.14 \\
Sparse & 0.6111 & 0.30 \\
Nonuniform & 0.7742 & 0.36 \\
\bottomrule
\end{tabular}
\caption{Optimal concurrent flow factors for three traffic matrices.}
\end{table}

\subsubsection{Uniform Traffic Matrix}

For the uniform traffic matrix, every node sends an equal amount of traffic to every other node. In this case, each non-diagonal entry of the traffic matrix is $4/7$.

Gurobi found an optimal concurrent flow factor of:

\[
\lambda = 0.7
\]

This means that the optimal topology can support 70\% of each requested traffic demand simultaneously.

\subsubsection{Sparse Traffic Matrix}

For the sparse traffic matrix, each node sends all 4 units of its traffic to a single destination. Specifically, node 0 sends to node 1, node 1 sends to node 2, and so on, with node 7 sending to node 0.

Gurobi found an optimal concurrent flow factor of:

\[
\lambda \approx 0.6111
\]

This means that each source-destination pair can send approximately:

\[
4(0.6111) \approx 2.4444
\]

units of traffic.

This result is lower than the uniform case. Since each source sends all its traffic to a single destination, some of the traffic must travel through intermediate nodes, which also use link capacity.

\subsubsection{Nonuniform Traffic Matrix}

For the nonuniform traffic matrix, each node distributes its 4 units of traffic across three destinations, with demands of 2, 1.5, and 0.5 units.

Gurobi found an optimal concurrent flow factor of:

\[
\lambda \approx 0.7742
\]

This was the highest concurrent flow factor among the three test cases. In this case, the traffic is distributed across multiple destinations, allowing the chosen topology to route a greater fraction of the requested traffic.

However, this does not necessarily mean that nonuniform traffic matrices will always have a higher concurrent flow factor. The result depends on the specific distribution of traffic.

\subsubsection{Overall Results}

For all three tests, Gurobi reported an optimality gap of 0\%, meaning that it verified the solutions were optimal under our formulation.

In addition, each resulting topology contained 32 directed links, with exactly 4 incoming and 4 outgoing links per node.

These results show that the optimal concurrent flow factor depends on how traffic is distributed across the network. They also demonstrate that our formulation can adapt the network topology to different traffic matrices rather than using the same fixed topology for every case.

\subsection{Conclusion}

Overall, our linear integer programming formulation allows us to find the optimal network topology for any given nonzero traffic matrix satisfying the hose model constraints, under our simple directed graph assumption.

By optimizing both the existence of links and the routing of traffic, we can maximize the concurrent flow factor while ensuring that all degree, capacity, and flow conservation constraints are satisfied.

Our Gurobi implementation successfully found optimal solutions for all three test matrices, demonstrating that the formulation works for different distributions of traffic.




\end{document}
